\input{./._relpath-to-latexroot.ltx}
\documentclass[\latexroot/\projectname]{subfiles}
\input{\latexroot/@resources/subfile-setup.ltx}
\begin{document}

\section{Conclusion}\whenintegrated{\label{conclusion}}
\whenintegrated{\label{sec:conclusion}} % integrated doc: SST for labels


The response of private consumer spending to a fiscal stimulus injection is at the heart of the income multiplier debate. We estimate a positive response of consumer spending to the Recovery

Act (2009-12) using regional variation. Localities that received $\$ 1$ more in government spending spent $\$ 0.29$ on non-durable spending and $\$ 0.09$ in auto purchases.

We estimate the aggregate response of consumer spending to fiscal stimulus using a quantitative equilibrium model. Our model is novel in that it embeds a regional framework into a heterogeneous agents, incomplete markets, new Keynesian model. The model generates a positive local non-durable consumption multiplier consistent with the data. This is a new finding and distinguishes our incomplete markets model from the existing literature that employed regional models with complete markets. The model predicts an aggregate non-durable consumption multiplier equal to 0.41 . This falls in the upper bound of estimates found in the literature (Hall, 2009).

Acknowledgments. We thank Dirk Krueger and five anonymous referees for extremely valuable advice. We also thank our discussants Karel Mertens and Claudia Sahm, as well as Gabriel Chodorow-Reich, Olivier Coibion, Erik Hurst, Greg Kaplan, Pete Klenow, David Lagakos, Alisdair McKay, Emi Nakamura, Valerie Ramey, Morten Ravn, Jon Steinsson, Gianluca Violante, and seminar participants at USCD, Wisconsin-Madison, UCL, Federal Reserve Bank of Minneapolis, Columbia GSB, NBER conference on Micro Data for Macro Models, and Barcelona Summer Forum 2018. The views expressed here are those of the authors and do not necessarily represent the views of the Federal Reserve Bank of St. Louis, the Federal Reserve Bank of Richmond, the Federal Reserve Bank of San Francisco, or the Federal Reserve System.

\section*{Supplementary Data}
Supplementary data are available at Review of Economic Studies online.

\section*{Data Availability Statement}
The non-confidential data and codes underlying this article are available in Zenodo, at \href{https://doi.org/}{https://doi.org/} 10.5281/zenodo.7483326. The confidential data underlying this article were provided by the Kilts Center for Marketing Data Center at The University of Chicago Booth School of Business (the Nielsen Dataset) and the New York Fed (Consumer Credit Panel-Equifax) under license.


% Smart bibliography: Only include bibliography if standalone AND has citations
\smartbib

\pdfonly{
% execute this version if compiling with pdflatex
  \captionsetup[figure]{list=no}
  \captionsetup[table]{list=no}
\end{document}
}
